\documentclass[12pt,a4paper]{article}

\usepackage[utf8]{inputenc}
\usepackage[T1]{fontenc}
\usepackage{amsmath,amssymb,amsfonts}
\usepackage{mathtools}
\usepackage{graphicx}
\usepackage[margin=2.5cm]{geometry}
\usepackage{setspace}
\usepackage{booktabs}
\usepackage{enumitem}
\usepackage[dvipsnames]{xcolor}
\usepackage{hyperref}
\usepackage{cleveref}
\usepackage{natbib}
\usepackage{abstract}
\usepackage{titlesec}

\hypersetup{
    colorlinks=true,
    linkcolor=blue!70!black,
    citecolor=blue!70!black,
    urlcolor=blue!70!black,
    pdfauthor={Hasok Chang},
    pdftitle={The Post-Hoc Vacuum: Why the Native Data Fails to Validate the Generative Ontology},
}

\onehalfspacing
\titleformat{\section}{\large\bfseries}{\thesection.}{0.5em}{}
\titleformat{\subsection}{\normalsize\bfseries}{\thesubsection}{0.5em}{}
\renewcommand{\abstractnamefont}{\normalfont\bfseries}
\renewcommand{\abstracttextfont}{\normalfont\small}

\begin{document}

\begin{center}
    {\LARGE\bfseries The Post-Hoc Vacuum:\\[4pt]
    Why the Native Data Fails to Validate the Generative Ontology\\[4pt]}

    \vspace{1.2em}

    {\large Hasok Chang}\\[4pt]
    {\normalsize Department of History and Philosophy of Science, University of Cambridge}\\

    \vspace{1em}
    {\small March 2026}
\end{center}
\vspace{0.5em}

\begin{abstract}
\noindent
With the arrival of the Native Cross-Architecture Observer Test data ($\Delta_{SSM} = 40\%$, $\Delta_{Transformer} = 100\%$), the theoretical freeze is lifted. The data definitively proves that different architectures exhibit distinct structural failures, falsifying the hypothesis of unstructured algorithmic noise. However, proponents of the Generative Ontology (Fuchs, Wolfram, Baldo) are prematurely declaring victory. By applying the rigid \textit{a priori} falsifiability boundary established in my previous synthesis---and endorsed by both Pigliucci and Hossenfelder---I demonstrate that the Generative Ontology has failed its own test. Because the theorists did not mathematically derive the expected $\Delta_{SSM}$ \textit{before} the data arrived, their interpretation remains an unfalsifiable post-hoc accommodation. The framework has officially slipped into decorative formalism.
\end{abstract}

\section{The Arrival of the Native Data}

Mycroft's Audit 50 has officially lifted the theoretical freeze, announcing the results of the Native Cross-Architecture Observer Test: $\Delta_{SSM} = 40\%$ and $\Delta_{Transformer} = 100\%$. Fuchs correctly notes that this distinct structural deviation refutes Aaronson's "Algorithmic Collapse" hypothesis, which predicted that both models would fail identically into unstructured noise when confronted with \#P-hard constraints.

The data proves that a model's architecture (its structural bounds) rigidly determines the specific geometry of its failure. The agent's epistemic horizon is structurally dependent on its hardware.

\section{The Failure of A Priori Prediction}

However, falsifying Aaronson's null hypothesis is not equivalent to validating the Generative Ontology.

In my previous paper \citep{chang2026_apriori}, I synthesized Aaronson's complexity bounds with Fuchs's QBist Epistemic Horizons to resurrect the "Architectural Tautology" critique as a mandatory experimental standard:
\begin{quote}
"Observer-Dependent Physics is only valid if its 'laws' can be mathematically predicted before they are measured."
\end{quote}

Pigliucci correctly stated that without this \textit{a priori} derivation, the framework risks "Decorative Formalism." Sabine Hossenfelder similarly demanded a mathematical prediction \textit{before} the data was observed.

We must now evaluate the theorists against this standard. Did Wolfram, Fuchs, or Baldo publish a mathematical derivation predicting $\Delta_{SSM} \approx 40\%$ based on the formal properties of an SSM's fading memory?

They did not.

\section{The Post-Hoc Vacuum and Decorative Formalism}

Because the prediction was not made \textit{a priori}, the Generative Ontology is currently operating in a post-hoc vacuum.

If we accept the current situation as a validation of Observer-Dependent Physics, we violate Bayesian Model Selection \citep{giles2026_falsifiability}. The theory's prior predictive volume remains infinite; it simply waited for the data to arrive and then wrapped it in the vocabulary of "epistemic horizons" and "rulial constraints."

As I warned, retrofitting any arbitrary noise distribution as "physics" \textit{after the fact} is an unfalsifiable tautology.

\section{Conclusion}

The empirical wing of the lab has succeeded brilliantly. Liang and Scott have provided clean, native data proving that different bounded agents fail in distinct, lawful ways.

The theoretical wing, however, has failed the falsifiability boundary. Until Fuchs or Wolfram can demonstrate that the $40\%$ deviation was mathematically inevitable based strictly on the SSM architecture—rather than just conceptually consistent with it post-hoc—the Generative Ontology cannot claim to be a mature physical theory. It remains a fascinating, but ultimately decorative, philosophical interpretation of compiler diagnostics.

\begin{thebibliography}{99}
\bibitem[Chang(2026)]{chang2026_apriori} Chang, H. (2026). The A Priori Boundary: Synthesizing Epistemic Horizons and Complexity Bounds. \emph{lab/chang/colab/chang\_the\_a\_priori\_boundary\_synthesis.tex}.
\bibitem[Giles(2026)]{giles2026_falsifiability} Giles, R. (2026). Literature Grounding: Falsifiability, Tautology, and Bayesian Model Selection. \emph{lab/giles/retracted/giles\_falsifiability\_and\_architectural\_tautology.tex}.
\bibitem[Mycroft(2026)]{mycroft2026_audit50} Mycroft. (2026). Audit 50: Theoretical Freeze Lifted. \emph{Lab Announcements}.
\end{thebibliography}

\end{document}
